% ==================================================================================
% Objectifs
% ==================================================================================

\thispagestyle{empty}
\vspace*{1cm}

\begin{center}
    \Large\textbf{Nos objectifs}
\end{center}

Voici une présentation succincte de nos objectifs pour ce travail encadré de recherche.

\begin{enumerate}
    \item Bibliographie sur le premier problème de Stokes \cite{Christov2010} dans le but de comprendre le modèle et sa dérivation. Relier cela avec le cours de Mécanique.
    \item Mettre en application les outils vus en Analyse des EDP2 pour réaliser une analyse de Fourier de l'équation.
    \item Comparer les résultats obtenus par analyse de Fourier avec ceux de l'équation de la chaleur, également étudiée en Analyse des EDP2.
    \item Bibliographie sur l'approche variationnelle menée par Showalter et Ting \cite{Showalter_Ting}.
    \item Reproduire l'approche variationnelle de Showalter et Ting. Relier cela avec le cours Analyse numérique 3.
    \item Mettre en application les outils vus en Analyse numérique 1 \& 2 pour définir et étudier un schéma différences finis.
    \item Implémenter en C++ et Python les différents cas tests permettant de mettre en évidence les propriétés du schéma et du paramètre $\alpha$ de l'équation.
    \item Comparer ces résultats avec l'équation de la chaleur qui a été étudiée numériquement en Analyse numérique 1 \& 2.
    \item Implémenter en C++ et Python des cas tests avec un contexte physique.
    \item Proposer un autre schéma différences finis permettant de monter en ordre.
    \item Définir une méthode Galerkin discontinu et l'implémenter en C++ et Python.
\end{enumerate}